\documentclass[11pt,a4paper]{article}
\usepackage[margin=25mm]{geometry}
\usepackage{amsmath,amssymb,amsthm,hyperref}
\setlength{\emergencystretch}{2em}
\newtheorem{theorem}{Theorem}
\newcommand{\gen}[1]{\langle#1\rangle}
\title{A singleton end forces a further unique crossing}
\author{Human Compression project: L044 continuation}
\date{}
\begin{document}\maketitle
Let $G=\gen{\rho_0,\ldots,\rho_{r-1}}\le A_n$, $n\ge12$,
be transitive with a connected string C-marking. Assume every
maximal ambient omission $G_j$ is intransitive and retain the
smaller-degree alternating induction of the complete critical-tail
reduction. Suppose $G_0$ has orbits $\{a\}$ and
$B=\Omega\setminus\{a\}$. Put $b=a^{\rho_0}$ and
$H=G_{\ge2}$. An eligible edge always crosses \emph{ambient}
$G_j$-orbits.

\begin{theorem}
If every colour $j\ge2$ has at least two eligible edges, then
$r\le\lfloor(n-1)/2\rfloor$. Consequently any counterexample
with this singleton end has a further colour $j\ge2$ with exactly
one ambient-eligible edge in the whole permutation graph.
\end{theorem}
\begin{proof}
Suppose the rank bound fails. The accepted critical odd-tail
reduction gives $n=2r\ge14$, $G_0=A_B$ naturally, and $H$
intransitive on its support. The group $H$ fixes $a,b$:
it fixes $a$ and commutes with $\rho_0$.
Write $f=|\operatorname{Fix}(H)|$ and let $c$ count its
nontrivial orbits. The full ambient interval theorem, using the
two eligible edges of every colour in the actual interval
$\{2,\ldots,r-1\}$, gives
\[
       2(r-2)\le n-f-c,\qquad f+c\le4.
\]
Since $f\ge2$ and $c\ge2$, necessarily $f=c=2$.
Thus $X=\Omega\setminus\{a,b\}$ is the union of exactly
two nontrivial $H$-orbits.

Write $\rho_0=(a\ b)\delta$, where $\delta$ is an odd
involution on $X$ centralizing $H$. If $\delta$ has no fixed
point on $X$, the complete accepted endpoint-centralizer theorem
already gives the rank bound. Its hypotheses are exactly the
present transitive even string C-action with intransitive maximal
omissions and singleton $G_0$-orbit. That theorem does not assume
that $H$ is transitive on $X$; its full actual-orbit and paired-row
arguments remain paid here.

We may therefore assume that $\delta$ has a fixed point in $X$.
Commutation permutes actual $H$-orbits. It cannot exchange the two
orbits, since that would make it fixed-point-free. On either
preserved orbit a centralizing permutation fixing one point fixes
every point. Thus $\delta$ is the identity on one orbit and acts
freely on the other; it cannot be the identity on both because
it is odd. Denote these orbits by $V$ and $O$, respectively.

For any colour $j\ge2$, an eligible $j$-edge in $O$ is carried
by $\rho_0$ to another eligible $j$-edge in $O$: the generators
commute, and $\rho_0\in G_j$ preserves each ambient omission
orbit. These two edges are distinct. Indeed a fixed unordered
edge of the fixed-point-free involution $\delta$ would itself
be a $\delta$-pair, and its endpoints would already be joined
by $\rho_0\in G_j$, contradicting eligibility.
If there is no eligible $j$-edge in $O$, its two assumed eligible
edges both lie in $V$, since $H$ fixes $a,b$.
Every colour therefore has an eligible pair in one actual
$H$-orbit. The complete paired-head theorem applies and gives
the contradiction.

Finally every colour has at least one eligible edge: if all its
edges stayed within omission orbits, adjoining its generator
could not turn the intransitive omission into transitive $G$.
The negation of the theorem's paired hypothesis consequently gives
exactly one eligible edge for some $j\ge2$, as stated.
\end{proof}

\paragraph{Scope and attribution.}
This is a new routing consequence for the singleton continuation of
Cameron--Fernandes--Leemans--Mixer, \emph{Highest rank of a polytope
for $A_n$}, arXiv:1605.09173v1, Propositions 5.31--5.32,
physical pages 42--45. It retains the complete critical-tail,
ambient-interval, endpoint-centralizer and paired-head proofs,
including their classical inputs and smaller-degree induction.
The new step converts global pairs to a pair in one actual orbit
using the exact two-orbit count and the fixed-point-free complement
theorem. No projected IP or local-to-ambient eligibility is inferred.
The further unique crossing and other full-source obligations still
require resolution. No new computation, formalization, original or
question credit is claimed.
\end{document}
